% ============================================================
\section{方法设计}
% ============================================================
\sectioncover{02}{方法设计}{}

% ====== 方法设计：提示词构建 ======

\begin{frame}{物理事件提示词构建}
  \slidekey{提示词构建是方法的一部分：把普通场景描述转化为可生成、可观察、可评价的物理事件。}

  \centering
  \includegraphics[width=0.88\paperwidth,height=0.55\paperheight,keepaspectratio]{prompt_dataset_pipeline.pdf}
\end{frame}

\begin{frame}{提示词构建：三阶段流程与设计原则}
  \slidekey{三阶段流程保证提示词质量，设计原则确保奖励聚焦物理事件。}

  \begin{columns}[T,onlytextwidth]
    \begin{column}{0.48\textwidth}
      \begin{block}{三阶段流程}
        \begin{compactenum}
          \item \textbf{候选生成：}围绕物理常识，由大语言模型批量构造描述。
          \item \textbf{物理筛选：}由 GPT-4o 判断候选是否真正涉及物理事件，剔除静态场景。
          \item \textbf{人工筛查：}处理边界样本与歧义结果，保证最终提示词可观察、可评价。
        \end{compactenum}
      \end{block}
    \end{column}%
    \begin{column}{0.48\textwidth}
      \begin{exampleblock}{设计原则}
        \begin{compactitem}
          \item \textbf{明确三要素：}每条提示词必须包含初始状态、动作过程和可观察后果。
          \item \textbf{排除不可评样本：}排除静态氛围、隐藏状态和多义结果。
          \item \textbf{聚焦物理事件：}让奖励聚焦物理事件本身，而非一般视觉质量。
        \end{compactitem}
      \end{exampleblock}

    \end{column}
  \end{columns}
\end{frame}

\begin{frame}{提示词生成与筛选模板}
  \slidekey{生成模板关注"初始状态、动作过程、可观察后果"，筛选模板剔除静态、不可评价或结果歧义样本。}

  \begin{columns}[T,onlytextwidth]
    \begin{column}{0.48\textwidth}
      \textbf{\textcolor{ZJUBlue}{生成候选提示词的模板}}
      \vspace{0.10cm}

      \begin{tikzpicture}
        \node[
          fill=CodeBg,
          draw=LineGray,
          rounded corners=3pt,
          text width=0.96\textwidth,
          align=left,
          inner sep=0.22cm,
          font=\footnotesize\ttfamily
        ] {
          Generate prompts testing\\
          physical commonsense.\\[3pt]
          Each prompt should:\\
          - Describe initial state,\\
          \phantom{- }action, consequence\\
          - Involve gravity, collision,\\
          \phantom{- }material, fluid, human\\
          \phantom{- }action, causal chains\\
          - Be observable in short\\
          \phantom{- }video, no audio/text
        };
      \end{tikzpicture}
    \end{column}%
    \begin{column}{0.48\textwidth}
      \textbf{\textcolor{ZJUBlue}{GPT-4o 筛选的提示词}}
      \vspace{0.10cm}

      \begin{tikzpicture}
        \node[
          fill=CodeBg,
          draw=LineGray,
          rounded corners=3pt,
          text width=0.96\textwidth,
          align=left,
          inner sep=0.22cm,
          font=\footnotesize\ttfamily
        ] {
          Evaluate prompt for\\
          physical commonsense test.\\[3pt]
          1. Truly involve physical\\
          \phantom{1. }event, not static\\
          2. Clear action and observable\\
          \phantom{2. }consequence\\
          3. Evaluable without audio,\\
          \phantom{3. }text, hidden states\\
          4. Unambiguous outcome\\[3pt]
          Score 1-5, where 5 = clear physical event
        };
      \end{tikzpicture}
    \end{column}
  \end{columns}
\end{frame}

% ====== 方法设计：后训练策略 ======

\begin{frame}{强化学习后训练策略总览}
  \slidekey{系统采用 Actor-VLM Judge-Trainer 架构：在线生成候选视频，视觉语言模型打分，强化学习更新 LoRA 参数。}

  \centering
  \includegraphics[width=0.90\paperwidth,height=0.59\paperheight,keepaspectratio]{vlm_feedback_physical_alignment.pdf}
\end{frame}

\begin{frame}{视觉语言模型双维度奖励}
  \slidekey{PC 奖励约束物理常识，SA 奖励约束语义完成度，两者共同减少奖励作弊。}

  \begin{columns}[T,onlytextwidth]
    \begin{column}{0.48\textwidth}
      \begin{block}{物理常识奖励 PC}
        \begin{compactitem}
          \item 重力方向与运动趋势是否合理。
          \item 碰撞响应、不可穿透约束。
          \item 材料形变、液体边界。
        \end{compactitem}
        \vspace{0.08cm}
        \centering \tagbox[ZJURed]{$r_{\mathrm{PC}}$}
      \end{block}
      \vspace{0.10cm}
      \begin{tikzpicture}
        \node[
          fill=CodeBg,
          draw=LineGray,
          rounded corners=3pt,
          text width=0.95\textwidth,
          align=left,
          inner sep=0.14cm,
          font=\scriptsize\ttfamily
        ] {
          Role: Expert in physical commonsense.\\
          Criteria: motion, force, gravity, collisions,\\
          \phantom{Criteria: }material, causality.\\
          Output: Score: X (1-5)
        };
      \end{tikzpicture}
    \end{column}%
    \begin{column}{0.48\textwidth}
      \begin{exampleblock}{语义贴合奖励 SA}
        \begin{compactitem}
          \item 主体、动作、场景是否出现。
          \item 关键视觉细节是否匹配。
          \item 防止静止化或模糊化逃避物理错误。
        \end{compactitem}
        \vspace{0.08cm}
        \centering \tagbox[ZJUGreen]{$r_{\mathrm{SA}}$}
      \end{exampleblock}
      \vspace{0.10cm}
      \begin{tikzpicture}
        \node[
          fill=CodeBg,
          draw=LineGray,
          rounded corners=3pt,
          text width=0.95\textwidth,
          align=left,
          inner sep=0.14cm,
          font=\scriptsize\ttfamily
        ] {
          Role: Expert video analysis judge.\\
          Criteria: objects, actions, scene,\\
          \phantom{Criteria: }temporal order, visual details.\\
          Output: Score: X (1-5)
        };
      \end{tikzpicture}
    \end{column}
  \end{columns}

\end{frame}

\begin{frame}{强化学习后训练策略优化}
  \slidekey{核心思想：用 VLM 奖励信号指导视频生成模型，通过截断策略梯度和 KL 正则化稳定训练。}

  \begin{block}{Flow-GRPO 的损失函数}
    \small
    \[
    L = \underbrace{-\mathbb{E}\!\left[\min\!\left(\rho_t\,\hat A,\,
    \operatorname{clip}(\rho_t,1\!-\!\varepsilon,1\!+\!\varepsilon)\,\hat A\right)\right]}_{L_{\mathrm{pg}}\text{：策略梯度损失}}
    \;+\;
    \underbrace{\beta \cdot L_{\mathrm{KL}}}_{\text{KL 正则项}}
    \]
  \end{block}

  \begin{columns}[T,onlytextwidth]
    \begin{column}{0.32\textwidth}
      \begin{exampleblock}{策略比值 $\rho_t$}
        \begin{minipage}[t][3.0cm]{\textwidth}
          \scriptsize
          \[
          \rho_t=\frac{p_\theta(z_{t-1}|z_t,c)}
          {p_{\theta_{\mathrm{old}}}(z_{t-1}|z_t,c)}
          \]
          {\normalfont 衡量新旧策略在同一转移上的概率比。$\rho_t>1$ 表示新策略更倾向于该动作。}
          \vfill
        \end{minipage}
      \end{exampleblock}
    \end{column}%
    \begin{column}{0.32\textwidth}
      \begin{exampleblock}{优势估计 $\hat A$}
        \begin{minipage}[t][3.0cm]{\textwidth}
          \scriptsize
          \[
          \hat A=w_{\mathrm{PC}}\hat A_{\mathrm{PC}}+
          w_{\mathrm{SA}}\hat A_{\mathrm{SA}}
          \]
          {\normalfont 群组内标准化后的奖励，$\hat A>0$ 表示优于平均水平。}\\[2pt]
          \[
          \hat A_h=\frac{r_h-\mu_{c,h}}{\sigma_{c,h}+\epsilon}
          \]
          \vfill
        \end{minipage}
      \end{exampleblock}
    \end{column}%
    \begin{column}{0.32\textwidth}
      \begin{exampleblock}{KL 正则项 $L_{\mathrm{KL}}$}
        \begin{minipage}[t][3.0cm]{\textwidth}
          \scriptsize
          \[
          L_{\mathrm{KL}}=\frac{\|\mu_\theta-\mu_{\mathrm{ref}}\|_2^2}{2\sigma_t^2\Delta t}
          \]
          {\normalfont 约束当前策略不偏离参考模型太远，防止灾难性遗忘。}
          \vfill
        \end{minipage}
      \end{exampleblock}
    \end{column}
  \end{columns}

  \vspace{0.06cm}
  \begin{tikzpicture}
    \node[
      fill=TitleBlue,
      draw=ZJUBlue!35,
      rounded corners=3pt,
      text width=0.96\textwidth,
      align=left,
      inner sep=0.14cm,
      font=\scriptsize
    ] {\textbf{直觉理解：}$\rho_t\hat A>0$ 时增大该转移概率，$\rho_t\hat A<0$ 时减小；clip 限制单步更新幅度，$\beta$ 控制保守程度。};
  \end{tikzpicture}
\end{frame}

\begin{frame}{代码实现流程}
  \slidekey{基于视觉语言模型评估的强化学习后训练策略的代码实现流程。}

  \small
  \textbf{输入：}基础生成器 $\pi_{\theta_0}$，参考策略 $\pi_{\mathrm{ref}}$，提示词集 $\mathcal{D}$，评估器 $J_{\mathrm{PC}}$ 和 $J_{\mathrm{SA}}$，组大小 $G$，KL 权重 $\beta$，截断范围 $\varepsilon$

  \vspace{0.12cm}
  \begin{tikzpicture}
    \node[
      fill=CodeBg,
      draw=LineGray,
      rounded corners=4pt,
      text width=0.96\textwidth,
      align=left,
      inner sep=0.20cm,
      font=\footnotesize
    ] {
      \begin{minipage}{\dimexpr\textwidth-0.44cm}
        \begin{tabbing}
        \hspace{1.6em}\= \hspace{1.6em}\= \hspace{1.6em}\= \hspace{1.6em}\= \kill
        初始化 $\pi_\theta\!\leftarrow\!\pi_{\theta_0}$，加载 LoRA 适配器 \hfill \textcolor{ZJUGray}{// 加载基础模型}\\
        \textbf{For} 每个训练轮次 \textbf{do}\\
        \> 从 $\mathcal{D}$ 采样提示词批次 $\mathcal{C}$ \hfill \textcolor{ZJUGray}{// 采样提示词}\\
        \> \textbf{For} 每个提示词 $c\in\mathcal{C}$ \textbf{do}\\
        \>\> \textbf{For} $i=1,\ldots,G$ \textbf{do} \hfill \textcolor{ZJUGray}{// 采样 $G$ 条候选}\\
        \>\>\> $(x_i,\tau_i,\ell_i^{\mathrm{old}})\!\leftarrow\!\textsc{SdeSample}(\pi_\theta,c)$\\
        \>\>\> $r_{i,\mathrm{PC}}\!\leftarrow\!J_{\mathrm{PC}}(x_i,c)$；$r_{i,\mathrm{SA}}\!\leftarrow\!J_{\mathrm{SA}}(x_i,c)$ \hfill \textcolor{ZJUGray}{// VLM 评分}\\
        \>\> \textbf{end for}\\
        \>\> $\hat A_{i,h}\!\leftarrow\!(r_{i,h}\!-\!\mu_{c,h})/(\sigma_{c,h}\!+\!\epsilon_{\mathrm{std}})$ \hfill \textcolor{ZJUGray}{// 群组内标准化}\\
        \>\> $A_i\!\leftarrow\!\operatorname{clip}(w_{\mathrm{PC}}\hat A_{i,\mathrm{PC}}\!+\!w_{\mathrm{SA}}\hat A_{i,\mathrm{SA}},-A_{\max},A_{\max})$ \hfill \textcolor{ZJUGray}{// 加权组合}\\
        \>\> \textbf{For} 每个转移 $(z_t,z_{t-1},c)$ \textbf{do}\\
        \>\>\> $\rho_i\!\leftarrow\!\exp(\log p_\theta(z_{t-1}|z_t,c)\!-\!\ell_i^{\mathrm{old}})$ \hfill \textcolor{ZJUGray}{// 策略比值}\\
        \>\>\> $L_{\mathrm{pg}}\!\leftarrow\!-\operatorname{mean}\{\min(\rho_i A_i,\operatorname{clip}(\rho_i,1\!-\!\varepsilon,1\!+\!\varepsilon)A_i)\}$\\
        \>\>\> $L_{\mathrm{KL}}\!\leftarrow\!\textsc{StepKL}(\pi_\theta,\pi_{\mathrm{ref}};z_t,c)$\\
        \>\>\> 更新 LoRA：$L_{\mathrm{pg}}+\beta L_{\mathrm{KL}}$ \hfill \textcolor{ZJUGray}{// 梯度更新}\\
        \>\> \textbf{end for}\\
        \> \textbf{end for}\\
        \textbf{end for}
        \end{tabbing}
      \end{minipage}
    };
  \end{tikzpicture}
\end{frame}

\begin{frame}{系统架构}
  \slidekey{Actor 节点、VLM Judge 节点、Trainer 节点解耦，使视频采样、VLM 评分和参数更新可以流水化执行。}

  \begin{columns}[T,onlytextwidth]
    \begin{column}{0.52\textwidth}
      \centering
      \includegraphics[width=\textwidth,height=0.52\paperheight,keepaspectratio]{system_architecture.pdf}
    \end{column}%
    \begin{column}{0.44\textwidth}
      \small
      \begin{compactitem}
        \item \textbf{Actor 节点}接收物理事件提示词，采样同一提示词下的 $G$ 条候选视频，记录旧策略概率，将视频发送给 Judge。
        \item \textbf{VLM Judge 节点}对每条视频独立输出 PC 与 SA 双维度奖励分数，将评分结果发送给 Trainer。
        \item \textbf{Trainer 节点}在群组内标准化奖励、计算优势估计，通过 Flow-GRPO 更新 LoRA 参数，再将新权重同步回 Actor。
      \end{compactitem}
    \end{column}
  \end{columns}
\end{frame}
